\documentclass[a4paper,10pt,twocolumn,oneside]{article}

% Column spacing for two-column mode
\setlength{\columnsep}{10pt} 
% Column separator line thickness for two-column mode
\setlength{\columnseprule}{0pt} 

% Math packages for definitions, theorems, and symbols
\usepackage{amsthm} 
\usepackage{amssymb}
\usepackage{amsmath}

% Font configuration package
\usepackage{fontspec} 

% Color packages
\usepackage{color}
\usepackage[x11names]{xcolor}

% Package for displaying source code
\usepackage{listings} 

% Header and footer customization
\usepackage{fancyhdr} 

% Graphics and image handling
\usepackage{graphicx} 

% List formatting
\usepackage{enumerate}

% Title formatting
\usepackage{titlesec}

% Drawing package
\usepackage{tikz}

% Code patching package
\usepackage{xpatch}

% Chinese typesetting support for XeLaTeX
\usepackage[CheckSingle, CJKmath]{xeCJK}

% Replaced CJKulem with ulem for underlining support compatible with xeCJK
\usepackage{ulem}

% Font encoding
\usepackage[T1]{fontenc}

% Additional utility packages
\usepackage{courier}

% Page margin and layout settings
\topmargin=0pt
\headsep=5pt
\textheight=780pt
\voffset=-40pt
\textwidth=545pt
\marginparsep=0pt
\marginparwidth=0pt
\marginparpush=0pt
\oddsidemargin=0pt
\evensidemargin=0pt
\hoffset=-42pt

%%%%%%%%%%%%%%%%%%%%%%%%%%%%%
% Font Settings
%%%%%%%%%%%%%%%%%%%%%%%%%%%%%

% Set main English font
\setmainfont{Consolas} 

% Set monospace font (Changed from Monaco to Courier New for Windows compatibility)
\setmonofont{Courier New} 

% Set main Chinese font (Microsoft JhengHei is natively supported on Windows)
\setCJKmainfont{Microsoft JhengHei} 

% Set section numbering depth to level 3
\setcounter{secnumdepth}{3} 

%%%%%%%%%%%%%%%%%%%%%%%%%%%%%
% Listings (Code Block) Settings
%%%%%%%%%%%%%%%%%%%%%%%%%%%%%
\makeatletter
\lst@CCPutMacro\lst@ProcessOther {"2D}{\lst@ttfamily{-{}}{-{}}}
\@empty\z@\@empty
\makeatother

\lstset{
    % The language of the code
    language=C++,
    % The size of the fonts that are used for the code
    basicstyle=\footnotesize\ttfamily,
    % The size of the fonts that are used for the line-numbers
    numberstyle=\footnotesize,
    % The step between two line-numbers. If it's 1, each line will be numbered
    stepnumber=1,
    % How far the line-numbers are from the code
    numbersep=5pt,
    % Choose the background color
    backgroundcolor=\color{white},
    % Show spaces adding particular underscores
    showspaces=false,
    % Underline spaces within strings
    showstringspaces=false,
    % Show tabs within strings adding particular underscores
    showtabs=false,
    % Adds a frame around the code
    frame=false,
    % Sets default tabsize to 2 spaces
    tabsize=2,
    % Sets the caption-position to bottom
    captionpos=b,
    % Sets automatic line breaking
    breaklines=true,
    % Sets if automatic breaks should only happen at whitespace
    breakatwhitespace=false,
    % If you want to add a comment within your code
    escapeinside={\%*}{*)},
    % If you want to add more keywords to the set
    morekeywords={*},
    % Keyword styling
    keywordstyle=\bfseries\color{Blue1},
    % Comment styling
    commentstyle=\itshape\color{Red4},
    % String styling
    stringstyle=\itshape\color{Green4},
}

%%%%%%%%%%%%%%%%%%%%%%%%%%%%%
% Document Body
%%%%%%%%%%%%%%%%%%%%%%%%%%%%%
\begin{document}

% Apply fancy header/footer style
\pagestyle{fancy}

\fancyhead[L]{Gugu Gaga Codebook}
\fancyhead[R]{\thepage}
\renewcommand{\headrulewidth}{0.4pt}
\renewcommand{\contentsname}{Contents} 

\scriptsize
\tableofcontents

%%%%%%%%%%%%%%%%%%%%%%%%%%%%%

\section{Math}
\subsection{Fpow}
\lstinputlisting{math/fpow.cpp}

\subsection{Combination}
\lstinputlisting{math/combination.cpp}

\subsection{Extgcd}
\lstinputlisting{math/extgcd.cpp}

\section{Shortest Path}
\subsection{Dijkstra}
\lstinputlisting{shortest_path/dijkstra.cpp}

\subsection{Bellman-Ford}
\lstinputlisting{shortest_path/bellman_ford.cpp}

\subsection{Floyd-Warshall}
\lstinputlisting{shortest_path/floyd_warshall.cpp}

\section{DP}
\subsection{0-1 Knapsack}
\lstinputlisting{dp/01_knapsack.cpp}

\subsection{Unbounded Knapsack}
\lstinputlisting{dp/unbounded_knapsack.cpp}

\subsection{SOS DP}
\lstinputlisting{dp/sos_dp.cpp}

\subsection{O(lgn) LIS}
\lstinputlisting{dp/olgn_lis.cpp}

\section{Graph}
\subsection{Euler Path}
\lstinputlisting{graph/euler_path.cpp}

\section{Data Structure}
\subsection{Segment Tree}
\lstinputlisting{data_structure/segment_tree.cpp}

\subsection{Binary Indexed Tree}
\lstinputlisting{data_structure/bit.cpp}

\section{Others}
\subsection{Discretization}
\lstinputlisting{others/discretization.cpp}

%%%%%%%%%%%%%%%%%%%%%%%%%%%%%
% Extended DP Section (From Slides)
%%%%%%%%%%%%%%%%%%%%%%%%%%%%%

\section{Advanced Dynamic Programming}

\subsection{Bitmask DP}
位元壓縮 DP，泛指狀態是由二進制表示的 DP。通常可以表示集合，0 表示不取，1 表示取。假設有 $n$ 個東西要考慮，使用長度為 $n$ 的 bitstring 表示狀態，複雜度通常為 $O(2^n)$。看到題目範圍 $n \le 20$ 即可考慮此解法。

\textbf{搭電梯問題 (Elevator Problem)} \\
$n$ 個人排隊，每個人重 $w_i$，要搭一台負重 $x$ 的電梯，問最少要搭幾趟？利用 $O(2^n)$ 窮舉最後上電梯的人，減少了窮舉 $n!$ 排列的時間複雜度。

\begin{lstlisting}
#define pp pair<int, int>
#define ff first
#define ss second

pp dp[(1<<20)]; // State up to 2^20

void solve() {
    int n, x;
    cin >> n >> x;
    vector<int> w(n);
    for (auto &c : w) cin >> c;
    
    // Initialize: maximum n rides, so set n+1 as infinity
    for(int s = 0; s < (1<<n); ++s) dp[s].ff = n + 1;
    
    // Empty set weight is x (forces first item to start a new ride)
    dp[0].ss = x, dp[0].ff = 0;
    
    // Iterate through all state sets
    for(int s = 1; s < (1<<n); ++s) {
        for(int i = 0; i < n; ++i) {
            // Check if person i is in the current set s
            if(!(s & (1<<i))) continue;
            
            // ... Transition logic goes here ...
        }
    }
}
\end{lstlisting}

\textbf{Hamiltonian Path / TSP} \\
圖中有 $n$ 個點，詢問是否能走過所有點且每個點僅經過一次 ($O(n \cdot 2^n)$)。需要記錄經過了哪些點，以及停留在哪個點。

\begin{lstlisting}
// Initialize: each node can reach itself
for(int i = 0; i < n; ++i) {
    dp[(1<<i)][i] = true;
}

// Iterate through all subsets
for(int s = 1; s < (1<<n); ++s) {
    for(int j = 0; j < n; ++j) {
        // Continue if vertex j is not in state s
        if(s & (1<<j)) continue;
        
        // Iterate all endpoints in state s
        for(int i = 0; i < n; ++i) {
            // s must pass through i
            if(!(s & (1<<i))) continue;
            
            // If there is a directed edge from i to j
            if(edge[i][j]) {
                dp[s | (1<<j)][j] |= dp[s][i];
            }
        }
    }
}

bool ans = false;
for(int i = 0; i < n; ++i) ans |= dp[(1<<n)-1][i];
\end{lstlisting}

\subsection{Sum Over Subsets (SOS DP) \& 高維前綴和}
目標是解決求給定集合的所有子集和：
$$F(S)=\sum_{x\subseteq S}A(x)$$
直接窮舉子集複雜度為 $O(3^n)$，使用 SOS DP 則可優化至 $O(n \cdot 2^n)$。

\textbf{高維前綴和概念 (High-Dimensional Prefix Sum)} \\
SOS DP 的本質其實就是高維度的前綴和推廣：
\begin{enumerate}
    \item \textbf{一維前綴和}：沿著單一方向累加，\texttt{F[x] += F[x-1]}。
    \item \textbf{二維前綴和}：先固定 $Y$ 對 $X$ 做一維前綴和 (\texttt{F[x][y] += F[x-1][y]})，再固定 $X$ 對 $Y$ 做一維前綴和 (\texttt{F[x][y] += F[x][y-1]})。幾何上就是把一條條的前綴和（線）疊加成一個矩形（面）。
    \item \textbf{N 維前綴和 (SOS DP)}：在 bitmask 中，共有 $n$ 個維度，每個維度的大小只有 0 和 1。我們只需依序窮舉每一維 $i$，若當前狀態 $s$ 在該維度是 1 (即 \texttt{s \& (1 << i)} 為真)，就將長度較短的狀態（即該維度為 0 的狀態 $s \oplus (1 \ll i)$）累加到 $F[s]$ 中。
\end{enumerate}

\begin{lstlisting}
// Precompute state weights (single point initialization)
for(int s = 0; s < (1<<n); ++s) {
    F[s] = A[s];
}

// Iterate through each dimension (bit)
for(int i = 0; i < n; ++i) {
    // Iterate through all states (points in the n-dimensional space)
    for (int s = 0; s < (1<<n); ++s) {
        // If the i-th bit of state s is 1 (meaning it includes this dimension)
        if (s & (1 << i)) {
            // Accumulate subset sums: F[s] += F[s - 1 on the i-th dimension]
            F[s] += F[s ^ (1 << i)];
        }
    }
}
\end{lstlisting}

\subsection{Counting DP}
用於計算方法數的 DP，通常可以配合排列組合與排容原理。

\textbf{障礙物棋盤走路計數問題} \\
給定 $n \times m$ 棋盤與 $N$ 個障礙物，只能向下或右走，求起點到終點且不經過障礙物的方法數。
利用總和扣除不合法路徑的方法：
$$dp[D] = cnt(O,D) - \sum_{\{k\}} cnt(k,D) \cdot dp[k]$$
轉移順序：將點先按 $X$ 排序，再按 $Y$ 排序，確保 partial-ordering 正確。

\begin{lstlisting}
vector<pair<int, int>> A(N);
sort(A.begin(), A.end());

// Insert start point at index 0 and end point at back
A.insert(A.begin(), pair<int, int>{1, 1}); 
A.push_back(pair<int, int>{n, m});

vector<int> dp(N + 5);

for(int i = 1; i <= N + 1; ++i) {
    // Calculate total valid permutations without obstacles
    dp[i] = cnt(0, i); 
    auto [x, y] = A[i];
    
    // Subtract combinations that pass through other obstacles
    for(int j = 1; j < i; ++j) {
        auto [xx, yy] = A[j];
        
        // Ensure obstacle j is strictly within the path to obstacle i
        if(xx <= x && yy <= y) {
            dp[i] -= cnt(j, i) * dp[j];
        }
    }
}
cout << dp[N+1];
\end{lstlisting}

\subsection{Probabilistic DP}
將機率丟到 DP 上做轉移，或者計算期望值 (利用期望值的線性性質 $E[X+c \cdot Y] = E[X] + c \cdot E[Y]$)。

\textbf{骰子問題 (Dice Probabilities)} \\
擲 $n$ 次六面骰子，所有結果相加落在 $[a, b]$ 的機率。

\begin{lstlisting}
// dp[i][j] represents the probability of reaching sum j after i rolls
double dp[N+1][6*N+1];

// Initialize the first roll probability
for(int i = 1; i <= 6; ++i) {
    dp[1][i] = 1.0 / 6.0;
}

for(int i = 2; i <= n; ++i) {
    // Max sum after n rolls is 6*n (assuming n <= 100, max 600)
    for(int w = 1; w <= 600; ++w) {
        // Iterate through all 6 dice faces
        for(int j = 1; j <= 6; ++j) {
            if(w >= j) {
                dp[i][w] += dp[i-1][w-j];
            }
        }
        // Divide once per state to prevent precision errors
        dp[i][w] /= 6.0;
    }
}

double ans = 0.0;
for(int i = a; i <= b; ++i) {
    ans += dp[n][i];
}
\end{lstlisting}

\textbf{抓老鼠期望值問題 (Mice Probability)} \\
袋中有 $w$ 隻白鼠、$b$ 隻黑鼠。雙方輪流抓，抓到白鼠贏，每次抓完黑鼠有機率跑出一隻。

\begin{lstlisting}
using ld = long double;
ld dp[1005][1005];

// Initialize: when only white mice remain, princess wins 100%
for (int i = 1; i <= w; i++) dp[i][0] = 1;
// Initialize: when only black mice remain, princess wins 0%
for (int i = 1; i <= b; i++) dp[0][i] = 0;

for (int i = 1; i <= w; i++) {
    for (int j = 1; j <= b; j++) {
        auto &cur = dp[i][j];
        
        // Case 1: Princess catches a white mouse immediately
        cur += (ld)i / (i + j);
        
        // Probability that both princess and dragon catch a black mouse
        ld x = (ld)j / (i + j) * (j - 1) / (i + j - 1);
        
        // Case 2: A black mouse escapes after they both catch black mice
        if (j >= 3) {
            cur += x * (j - 2) / (i + j - 2) * dp[i][j-3];
        }
        
        // Case 3: A white mouse escapes after they both catch black mice
        if (j >= 2) {
            cur += x * i / (i + j - 2) * dp[i-1][j-2];
        }
    }
}
\end{lstlisting}

\end{document}